\subsection{Compiling the UI}
The graphical user interface of the Grimace Box recording application was designed using QT Designer.
The \verb!.ui! files produced by QT Designer need to be converted to Python files.
The newest versions of the python files are uploaded to the Github repository alongside the rest of the code,
so you only need to regenerate them if you have edited the \verb!.ui! files.

To do this, you can run the \verb!compile_ui.bat! batch script to convert all \verb!.ui! files
or use the \verb!pyside6-uic! command to convert individual files.
The \verb!.ui! files are located in the \verb!ui! folder,
the generated python files are placed in the \verb!app/layout! folder.

\subsection{Building the Executable}
To deploy the application to a client system, it needs to be compiled to an executable (\verb!.exe!) file.
This is done using \verb!pyinstaller!, which packs the python code together with all
required libraries and a Python interpreter inside an executable file.

The Python dependencies are managed with \verb!uv!.
Run \verb!uv sync! in the project folder to install the exact dependency versions from \verb!uv.lock!
into the \verb!.venv! folder.

To create the executable, run \verb!uv run cmd /c build_application.bat!.
A folder called \verb!dist/GrimAce! is created, which contains the executable file as well as a
folder called \verb!_internal!, containing the python interpreter, libraries, license files and additional files.
The executable file only works if it is in the same location as the \verb!_internal! folder.
If you want to move the executable file, you need to move the entire \verb!GrimAce! folder.

Releases are also built automatically by a GitHub Actions workflow when a release is published on GitHub.